% Copyright (c) 2026 Eric van der Merwe
% SPDX-License-Identifier: MIT

The Book of Revelation was written by a Christian prophet named John, believed to be John the Apostle, while exiled on the island of Patmos, around 90-95 AD (approximately 65 years after the death of Christ). Many scholars distinguish this John from the author of the Fourth Gospel, though this has not been settled. It is the New Testament's only fully apocalyptic book, combining short letters to seven churches in Asia Minor with a sequence of symbolic visions.

John was writing to congregations under pressure from the Roman imperial cult and from compromise within the churches themselves, and could therefore encode his consolation and warning in the language of the prophets rather than in narrative biography (for example, the beasts and the fall of Babylon draw on Daniel and Ezekiel, rather than naming Rome directly).

Those visions have been read in more than one way. Two accounts of the book's message are given here: the evangelical Protestant reading common among the churches of the Reformation's heirs, and the Swedenborgian reading that informs the notes in this edition.

\medskip
\noindent{\bf Evangelical Protestant interpretation.}

Evangelical Protestants take Revelation as inspired Scripture whose centre is Jesus Christ: the faithful witness, firstborn from the dead, and ruler of the kings of the earth, who was slain and is alive for evermore. Its theme is that history belongs to the Lamb. He alone is worthy to open the sealed scroll of God's purpose; the churches are called to repent of compromise and to endure; the powers that parody his rule (the dragon, the beasts, Babylon) are judged; and the story ends not in escape from creation but in its renewal, a new heaven and a new earth where God dwells with his people. Evangelicals do not agree on the timetable. Some read the visions as largely fulfilled in the first century (preterist), some as a symbolic map of the whole church age (idealist), some as a panorama of later history (historicist), and many, especially in popular evangelicalism, as still-future tribulation, Antichrist, and millennial reign (futurist). Those debates matter, but they sit inside a shared confession: Revelation is a pastoral prophecy, not a secret code detached from the gospel. Its practical word is to worship God, refuse the world's idolatry, and pray ``Come, Lord Jesus.''

\medskip
\noindent{\bf Swedenborgian interpretation.}

Swedenborgians read the same book as a spiritual, not a chronological-political, prophecy. In Emanuel Swedenborg's {\it Apocalypse Revealed} and {\it Apocalypse Explained}, Revelation describes the end of a corrupted church and the Lord's establishing of a new one, together with a last judgment that, they hold, was accomplished in the spiritual world rather than as the destruction of the physical universe. The imagery is correspondence throughout. The seven churches are states of life and doctrine among those who belong to the church; horses, seals, trumpets, and vials are successive uncoverings of interior evil and falsity; Babylon is religion turned into dominion; the dragon and the beasts are faith separated from charity. The Lamb is the Lord in his Divine Human; the holy city, New Jerusalem, is the New Church, in which faith and charity, good and truth, are conjoined as in a marriage. ``The time is at hand'' is therefore the Lord's making all things new in heaven and in the church, and the book's comfort is that this new Christianity is the Lord's own work, not a forecast of earthly catastrophe.
